\section{Attaining the constructive cutoff budget}\label{sec:attained-cutoff}
The preceding fixed-background argument uses strict slack in both the
floor and the tail budget. We now remove both restrictions. The cutoff
event remains strictly $r>K$; mass exactly at $K$ is essential.

\begin{theorem}\label{thm:attained-cutoff}
Fix $0<\rho<1$, $s=(1+\rho)/(1-\rho)$, and
\[
 k_s=\frac{s-1}{s}<K<1,\qquad
 g(K)=\frac{s^s}{(s-1)^{s-1}}(1-K)K^{s-1}.
\]
For every integer $L\ge1$ there is a strictly positive density $r_L$
on the uniform Boolean cube such that
\[
 \mathbb Er_L=1,\qquad r_L\ge K,\qquad
 \mathbb E[r_L^s\mathbf1_{r_L>K}]\le1+g(K),
\]
and $\lambda_1(T_{r_L})\to1$, $\lambda_2(T_{r_L})\to1$.
Consequently $\gap(T_{r_L})\to0$. The same family meets every
prescribed floor $0<b\le K$ and every budget $\tau\ge1+g(K)$.
This is an attained constructive upper budget, not a claim that
$1+g(K)$ is the optimal threshold over all collapsing families.
\end{theorem}

\begin{proof}
All expectations below use normalized uniform probability. Write
$Q_d=P_{\sqrt\rho,d}$. For positive $r$, the full spectra of
$T_r=M_{\sqrt r}P_{\rho,d}M_{\sqrt r}$ and $S_r=Q_dM_rQ_d$
agree, by the $AA^*/A^*A$ identity with $A=M_{\sqrt r}Q_d$.

\emph{The inherited balanced tilt.}
Use the explicit family in the proof of
Theorem~\ref{thm:critical-exact}. Fix a sufficiently small $d>0$ and set
\[
 \delta_n=dn^{-1/3},\qquad
 N(\delta)=\left(\frac{(1+\delta)^s+(1-\delta)^s}{2}\right)^{1/s},
 \qquad
 R_n(y)=\prod_{i=1}^n\left(\frac{1+\delta_n y_i}{N(\delta_n)}\right).
\]
In particular, the normalization is applied to every product factor.
Define
\[
 H_n=Q_nM_{R_n}Q_n,\qquad
 \kappa(\delta)=\frac{1+\rho+
 \sqrt{(1-\rho)^2+4\rho\delta^2}}{2N(\delta)},\qquad
 b_n=\kappa(\delta_n)^{-n},\qquad
 F_n=P_{\rho,n}+b_nH_n.
\]
The earlier proof gives
\begin{equation}\label{eq:attained-inherited}
 b_n\to1,\quad \|H_n\|=\kappa(\delta_n)^n\to1,\quad
 \lambda_1(F_n)\to1,\quad\lambda_2(F_n)\to1,
 \quad\mathbb ER_n^s=1.
\end{equation}
These are conclusions for that explicit family, not additional hypotheses.
For completeness, its interior moments have the exact expression
\[
 a_n(t):=\mathbb ER_n^t
 =\left(\frac{(1+\delta_n)^t+(1-\delta_n)^t}
 {2N(\delta_n)^t}\right)^n\longrightarrow0\qquad(0<t<s).
\]
Indeed their logarithm is
$n[t(t-s)\delta_n^2/2+O(\delta_n^4)]$.
At the critical exponent, $\log\kappa(\delta)=O(\delta^4)$.
The block Gram comparison proved earlier has two diagonal top
eigenvalues one and an off-diagonal norm tending to zero; it controls
the actual two leading eigenvalues of $F_n$. Its second lower bound
is eventually positive, so it is not an added zero of the Gram matrix.

\emph{Growing finite background laws.}
The scalar optimizer from Proposition~\ref{prop:background-optimum} has
\[
 v_*=\frac{sK}{s-1}>1,\qquad
 p_*=\frac{(s-1)(1-K)}K\in(0,1),\qquad p_*v_*^s=g(K).
\]
For sufficiently large $L$ choose
\[
 k=\lfloor L^{1/3}\rfloor,\quad
 M=\left\lfloor\frac{L}{2k}\right\rfloor,\quad n=L-kM,
 \qquad p=p_L=\frac{\lfloor2^kp_*\rfloor}{2^k},
 \qquad v=v_L=K+\frac{1-K}{p}.
\]
Then $p\in[p_*/2,p_*]$, $p\to p_*$, $v\to v_*$ and
$v\le V:=2v_*$. Also $k,M,n\to\infty$, $n\ge L/2$ and
$M/k\to\infty$. Choose exactly $2^kp$ patterns in $\{-1,1\}^k$,
and let $f_L$ take value $v$ there and $K$ elsewhere.
Replace its $k$ signs by the parities of $k$ disjoint blocks of $M$
bits. Call the resulting background $A_L$, with high-sector indicator
$I_L$. Independent unbiased parities give
\[
 \mathbb EI_L=p,\qquad A_L=K+(v-K)I_L,\qquad\mathbb EA_L=1.
\]
For a nonempty Walsh support $U$,
\[
 Q_dM_{\chi_U}Q_d\chi_S
 =\rho^{|U|/2+|S\setminus U|}\chi_{S\mathbin\triangle U}.
\]
Thus this weighted permutation has norm exactly $\rho^{|U|/2}$.
Each Walsh coefficient of $f_L$ has absolute value at most
$\mathbb E f_L=1$, and there are at most $2^k-1$ nonconstant ones.
Every substituted support has size at least $M$, so
\begin{equation}\label{eq:attained-parity}
 \|Q_{kM}M_{A_L}Q_{kM}-P_{\rho,kM}\|
 \le(2^k-1)\rho^{M/2}\longrightarrow0.
\end{equation}
The limit follows from $k\log2+(M/2)\log\rho\to-\infty$.
In particular, the growing number of coefficients is explicitly
controlled. The scalar affine minorant gives
\begin{equation}\label{eq:attained-dyadic-excess}
 d_L:=p_Lv_L^s-g(K)\ge0,\qquad d_L\to0.
\end{equation}
No uniform finite cube is assumed to realize a non-dyadic $p_*$ exactly.

\emph{Clipping, exact mean, and the amplitude reserve.}
Put $\theta=\min(1,s/2)$, so both $\theta$ and $s-\theta$ lie in $(0,s)$.
For $n=n_L$ define
\[
 c_n=a_n(\theta)^{1/(2\theta)},\qquad
 U_n=(R_n-c_n)_+,\qquad m'_n=\mathbb EU_n.
\]
Then $c_n\to0$, $m'_n\le a_n(1)\to0$, and
\begin{equation}\label{eq:attained-clipping}
 \mathbb P(U_n>0)\le\sqrt{a_n(\theta)},\qquad
 \|Q_nM_{U_n-R_n}Q_n\|\le c_n.
\end{equation}
The first inequality is Markov's inequality at $c_n$; the second uses
$0\le R_n-U_n\le c_n$ and $\|Q_n\|=1$.
Equality $R_n=c_n$ is in the complement of the strict event.

For real $s>1$, $c>0$ and $z\ge0$, the mixed-power bound is
\begin{equation}\label{eq:attained-mixed}
 0\le(c+z)^s-c^s-z^s
 \le B_s\bigl(c^{s-\theta}z^\theta+c^\theta z^{s-\theta}\bigr),
 \qquad B_s=s2^{s-1}.
\end{equation}
To see the upper bound, divide by $c^s$. For $z/c\le1$ the
mean value theorem bounds the difference by $B_s(z/c)$, which is
at most $B_s(z/c)^\theta$. For $z/c\ge1$ exchange $c,z$.
The lower bound follows by integrating a positive derivative.
Thus this estimate does not require an integer exponent.
Set
\begin{align*}
 E_L={}&K^s\sqrt{a_n(\theta)}\\
 &+B_s\bigl[(K^{s-\theta}+V^{s-\theta})a_n(\theta)
 +(K^\theta+V^\theta)a_n(s-\theta)\bigr],\\
 e_L={}&d_L+E_L.
\end{align*}
Both are nonnegative and tend to zero. Eventually $e_L<1/2$ and
$m'_n<1-K$. On that range let
\begin{equation}\label{eq:attained-density}
 \beta_L=(1-e_L)^{1/s},\qquad
 \gamma_L=\frac{\beta_Lm'_n}{p_L},\qquad
 r_L(x,y)=A_L(x)-\gamma_L I_L(x)+\beta_LU_n(y).
\end{equation}
Here $0<\beta_L\le1$, $\beta_L\to1$ and $\gamma_L\to0$.
The mean is exactly $1-p_L\gamma_L+\beta_Lm'_n=1$.
The low sector has density $K+\beta_LU_n\ge K$; on the high sector,
\[
 v_L-\gamma_L=K+\frac{1-K-\beta_Lm'_n}{p_L}>K.
\]
Thus mean compensation has not lowered the floor. It acts only
on the high sector, and no global normalization is applied.

On the low sector the strict event is precisely $U_n>0$. Apply
\eqref{eq:attained-mixed} with $c=K$ and $z=\beta_LU_n$, restricted
to that event. Its constant contribution is at most
$K^s\mathbb P(U_n>0)$. The high sector is entirely above $K$;
discarding the nonnegative subtraction $\gamma_L$ gives an upper
bound from \eqref{eq:attained-mixed} with $c=v_L$.
Averaging these bounds with probabilities $1-p_L,p_L$, the
coefficient of $\beta_L^s\mathbb EU_n^s$ is exactly one.
Use $U_n\le R_n$, $v_L\le V$, $\beta_L\le1$,
$\mathbb ER_n^s=1$, and \eqref{eq:attained-clipping}. The two mixed
terms are bounded by those defining $E_L$. Hence
\begin{equation}\label{eq:attained-exact-budget}
 \mathbb E[r_L^s\mathbf1_{r_L>K}]
 \le p_Lv_L^s+\beta_L^s+E_L
 =g(K)+d_L+1-e_L+E_L=1+g(K).
\end{equation}
The reserve pays for the dyadic excess and every retained tail error.
No convergence of $s$-moments was inferred from convergence in probability.

\emph{The actual leading eigenvalues.}
Write $P_b=P_{\rho,kM}$, $P_t=P_{\rho,n}$ and similarly $Q_b,Q_t$.
By linearity and tensorization,
\[
 S_{r_L}=(Q_bM_{A_L}Q_b-\gamma_LQ_bM_{I_L}Q_b)\otimes P_t
 +\beta_LP_b\otimes(Q_tM_{U_n}Q_t).
\]
Compare with $B_L=P_b\otimes F_n$. Since $0\le I_L\le1$,
\eqref{eq:attained-parity} and \eqref{eq:attained-clipping} imply
\begin{align*}
 \|S_{r_L}-B_L\|
 &\le(2^k-1)\rho^{M/2}+\gamma_L
       +\beta_Lc_n+|\beta_L-b_n|\|H_n\|\longrightarrow0.
\end{align*}
All four errors vanish by the established limits.
Both tensor factors of $B_L$ are positive definite and
$\lambda_1(P_b)=1$. Therefore
\[
 \lambda_1(B_L)=\lambda_1(F_n),\qquad
 \lambda_2(F_n)\le\lambda_2(B_L)\le\lambda_1(F_n).
\]
The lower bound uses two actual product eigenvectors with constant
background factor. In view of \eqref{eq:attained-inherited}, both
leading eigenvalues of $B_L$ tend to one. Weyl's self-adjoint
perturbation bound transfers both limits to $S_{r_L}$, and the
full-spectrum identity transfers them to $T_{r_L}$. This controls
the top eigenvalue as well as the second; two trial vectors alone
would not suffice.

Choose one $L_0$ so that every eventual condition above holds for
all $L\ge L_0$. For $1\le L<L_0$ set $r_L=1$. Its floor is at
least $K$, its mean is one, and its strict $K$-tail cost is
$1\le1+g(K)$. This supplies every dimension without changing the
limits. No effective estimate of $L_0$ is asserted.
\end{proof}

For the non-strict event $r_L\ge K$, every point contributes and the
tail becomes the total moment. The theorem therefore cannot be
relabeled with a non-strict event. The scalar optimization, clipping
and perturbation tools are classical; this construction and its
source audit do not establish external priority. The unresolved
universal interval between the sufficient region
$\tau<2-K^{s-1}$ and this constructive upper budget remains.
No marginal matching, locality, heat-bath gap or physical conclusion
is added to the separate families proved earlier.
